\documentclass[10pt,a4paper]{amsart}
\usepackage[margin=19mm]{geometry}
\usepackage{amsmath,amssymb,mathtools}
\usepackage[scr=rsfs,cal=euler]{mathalfa}
\usepackage{xfrac}
\usepackage[unicode,hidelinks,bookmarks=false]{hyperref}
\newcommand{\ie}{i.e.\ }
\newcommand{\cf}{cf.\ }
\newcommand{\resp}{respectively\ }
\newcommand{\Subsection}{\subsection}
\providecommand{\nobreakdash}{\nobreak\hbox{-}}
\setlength{\emergencystretch}{2em}
\multlinegap0pt
\usepackage{amssymb}
\usepackage{bm}
\usepackage[shortlabels]{enumitem}
\usepackage{mathtools}
\usepackage{xcolor}
\usepackage{float}
\newtheorem{theorem}{Theorem}
\newtheorem{lemma}{Lemma}
\def\red#1{\textcolor{red}{#1}}
\newcommand{\wtb}{\langle\partial_Z\rangle }
\title{The stability threshold for Boussinesq equations around stratified Couette flow with unequal viscosity and thermal diffusivity on $\mathbb{T}\times\mathbb{R}^2$}
\author{Pengyu Chen, Zeren Zhang}
\date{Source: arXiv:2609.29916v1 (CC BY 4.0)}
\begin{document}
\maketitle

\noindent\emph{Source and licence.} The stability threshold for Boussinesq equations around stratified Couette flow with unequal viscosity and thermal diffusivity on $\mathbb{T}\times\mathbb{R}^2$. Version arXiv:2609.29916v1 (CC BY 4.0), distributed by arXiv under the Creative Commons Attribution 4.0 International licence, http://creativecommons.org/licenses/by/4.0/, as recorded in the arXiv licence metadata for this exact version and shown on the abstract page https://arxiv.org/abs/2609.29916v1. Full licence notice: \texttt{licenses/arxiv-2609.29916-CC-BY-4.0.txt}.\par\smallskip

\noindent\textit{Editorial scope.} Counted unit: the article's lemma est-fun (source0.tex lines 419--431). The article's complete original proof of it is delivered with it (source0.tex lines 432--434, one \texttt{proof} environment of 3 lines). No mathematics is written, repaired or re-derived: the statement and the proof are the article's own lines, in the article's own notation, and no retained line is substituted or re-flowed.  Retained uncounted as the source-local context the proof needs: lines 88--97; lines 123--140; lines 380--385.  Nothing was omitted from any retained span, so the package carries no omission record.  No retained line contains a citation, so the wrapper carries no bibliography and no bibliography entry.  No retained line contains a figure environment, an image-inclusion command or drawing code, so no image is delivered and the package declares no asset dependency.  Editorial formatting only, in addition: an AMS \texttt{amsart} preamble that restates the article's own declarations and macros used by the retained lines, namely the environments \texttt{theorem}, \texttt{lemma} on separate counters, together with the standard packages those lines need.  The retained environments print in this document as Theorem 1, Lemma 1 in order of appearance, whereas the article prints the same items with the numbering of its own full document.  The complete native source of the article as distributed in the arXiv e-print for this exact version is bundled unchanged as \texttt{sources/211709/source0.tex}.
	\begin{equation}\label{eq-main}
		\left\{
		\begin{aligned}
			&\partial_t u+y\partial_x u+(u^2,0,0)^T+(u\cdot\nabla)u+\nabla p^L+\nabla p^{NL}=\kappa \Delta u-\beta(0,\theta,0)^T,\\
			&\partial_t \theta+y\partial_x \theta-\beta u^2+(u\cdot\nabla)\theta=\mu\Delta\theta,\\
			&\Delta p^L=-2\partial_x u^2-\beta\partial_y \theta,\quad\Delta p^{NL}=-\partial_i u^j\partial_j u^i,\\
			&(u,\theta)|_{t=0}=(u_{in},\theta_{in}),
		\end{aligned}
		\right.
	\end{equation}

	\begin{theorem}\label{mainthe}
		Assume that $\beta>0$, $0<\kappa,\mu\leq1$, and $C_1^{-1}\mu\leq\kappa\leq C_1\mu$ for some constant $C_1\geq1$. Set $\nu=\min\{\kappa,\mu\}$ and let $(u_{in},\theta_{in})$ be the initial datum for \eqref{eq-main}. For any $0<a_0\ll1$ and $s\geq4$, there exists $\epsilon_0=\epsilon_0(\beta,s,a_0,C_1)>0$ such that, if it holds that
		\begin{equation}
			\|\langle \nabla_{x,z}\rangle (u_{in},\theta_{in})\|_{H^{s+1}}+\|\wtb(u_{0,in},\theta_{0,in})\|_{W^{s,1}}=\epsilon\leq\epsilon_0\nu^{\frac{2}{3}}|\ln{\nu}|^{-2-2a_0},
		\end{equation}
		then the system \eqref{eq-main} admits a unique global solution. Moreover, the profile $(U,\Theta)(t,x,y,z)=(u,\theta)(t,x+yt,y,z)$ satisfies the following estimates:
		\begin{subequations}\label{est-main}
			\begin{align}
				&\|e^{a_1\nu^{\frac{1}{3}}t}\partial_x\nabla_{x,z}U^1_{\neq}\|_{L^\infty H^s}+\nu^{\frac{1}{6}}\|\partial_x\nabla_{x,z}U^1_{\neq}\|_{L^2H^s}\lesssim|\ln\nu|^{1+a_0}\epsilon,\label{main-u1}\\
				&\|e^{a_1\nu^{\frac{1}{3}}t}\nabla_{x,z}\nabla_LU^2_{\neq}\|_{L^\infty H^s}+\|\partial_x\nabla_{x,z}U^2_{\neq}\|_{L^2 H^s}+\nu^{\frac{1}{6}}\|\nabla_{x,z}\nabla_LU^2_{\neq}\|_{L^\infty H^s}\lesssim\epsilon,\label{main-u2}\\
				&\|e^{a_1\nu^{\frac{1}{3}}t}\nabla_{x,z}^2U^3_{\neq}\|_{L^\infty H^s}+\nu^{\frac{1}{6}}\|\nabla_{x,z}^2U^3_{\neq}\|_{L^2H^s}\lesssim|\ln\nu|^{1+a_0}\epsilon,\label{main-u3}\\
				&\|e^{a_1\nu^{\frac{1}{3}}t}\nabla_{x,z}^2\Theta_{\neq}\|_{L^\infty H^s}+\nu^{\frac{1}{6}}\|\nabla_{x,z}^2\Theta_{\neq}\|_{L^2H^s}\lesssim\epsilon,\label{main-theta}\\
				&\|\wtb U^1_0\|_{L^\infty H^s}+\nu^{\frac{1}{6}}\| U^1_0\|_{L^\infty H^{s+1}}+\|\wtb (U^{2}_0,U^3_0,\Theta_0)\|_{L^\infty H^{s+1}}\lesssim\epsilon,\label{main-u10}\\
				&\nu^{\frac{1}{12}+a_2}\|\wtb(U^2_0,\Theta_0)\|_{L^2W^{s-2,\infty}}+\nu^{a_2}\|\wtb U^3_{0}\|_{L^2W^{s-1}}\lesssim\epsilon.\label{main-u230}
			\end{align}
		\end{subequations}
		Here $a_1,a_2$ are sufficiently small constants, $\nabla_{x,z}=(\partial_x,\partial_z)$, and $\nabla_L=(\partial_x,\partial_y-t\partial_x,\partial_z)$.
	\end{theorem}

			\begin{align}
			&\max\left\{\langle t\rangle^{-2},\nu^{2/3}\right\}\lesssim m\leq 1;\label{prop1m}\\
			& k^2+l^2\lesssim p m;\label{prop2m}\\
			& m(k,\eta,l)\lesssim \langle\eta-\eta_1,l-l_1\rangle^2m(k,\eta_1,l_1),\label{prop3m} \\
            & \left|\frac{\partial_t p}{2p}\right| \leqslant -\frac{\partial_t m}{2m} + \frac{\nu}{16} p. 
		\end{align}

    \begin{lemma}\label{est-fun}
        Under the assumptions of Theorem \ref{mainthe}, it holds that
        \begin{align*}
		%&\red{\|m^{1/2} p^{\frac{1}{2}}\partial_X U^1_{\neq}\|_{L^\inftyH^s}+\nu^{1/6}\|m^{1/2} p^{\frac{1}{2}}\partial_X U^1_{\neq}\|_{L^2 H^s}\lesssim\epsilon |\ln{\nu} |^{1+a_0} ,}\\
        &\|q^{\frac{1}{2}}\partial_X U^1_{\neq}\|_{L^\infty H^s}+\nu^{\frac{1}{6}}\|q^{\frac{1}{2}}\partial_X U^1_{\neq}\|_{L^2 H^s}\lesssim\epsilon|\ln{\nu} |^{1+a_0},\\
		&\|(pq)^{\frac{1}{2}}U^2_{\neq}\|_{L^\infty H^s}+\|\partial_Xq^{\frac{1}{2}}U^2_{\neq}\|_{L^2 H^s}+\nu^{\frac{1}{6}}\|(pq)^{\frac{1}{2}}U^2_{\neq}\|_{L^2 H^s}\lesssim\epsilon,\\
		&\|q U^3_{\neq}\|_{L^\infty H^s}+\nu^{\frac{1}{6}}\|q U^3_{\neq}\|_{L^2 H^s}\lesssim\epsilon|\ln{\nu} |^{1+a_0},,\\
		&\|q\Theta_{\neq}\|_{L^\infty H^s}+\nu^{\frac{1}{6}}\|q\Theta_{\neq}\|_{L^2H^s}\lesssim\epsilon,\\
		&\|\wtb U^1_{0}\|_{L^\infty H^s}+\nu^{\frac{1}{6}}\|U^1_{0}\| _{L^\infty H^{s+1}}+\|\wtb (U^2_{0},U_0^3,\Theta_0) \|_{L^\infty H^{s+1}} \lesssim\epsilon.
        %    &\|U^1_{0}\|_{L^\infty H^{s+1}}\lesssim\epsilon \nu^{-1/6},\\
		%&\|(1+\partial_Z )(U^2_{0},U_0^3,\Theta_0) \|_{L^\infty H^{s+1}}+\nu^{1/2} \|(1+\partial_Z )(U^2_{0},U_0^3,\Theta_0) \|_{L^2 H^{s+1}} \lesssim\epsilon.
	\end{align*}
    \end{lemma}

    \begin{proof}
        Lemma \eqref{est-fun} follows from the bootstrap hypotheses, \eqref{prop2m} and the divergence free condition, we thus omit its proof. 
    \end{proof}

\end{document}
